\section{Kiểm thử và đánh giá}
\label{chap:eval}

Đánh giá gồm ba tiêu chí theo CO5173: tính đúng sau xử lý, hiệu năng và hiệu quả khai thác dữ liệu. Tính đúng được đối chiếu với oracle và trạng thái quan sát qua API. Hiệu năng dùng độ trễ từ mốc đầu đến kết quả đã lưu hoặc API và thông lượng theo số sự kiện duy nhất. Hiệu quả được kiểm tra bằng thao tác nghiệp vụ trong ứng dụng và bằng nhãn rút học ngoại tuyến của OULAD. Nhật ký gốc và cách chạy nằm ở Phụ lục~\ref{app:run}.

\subsection{Chiến lược kiểm thử}

\begin{table}[H]
\centering\small
\caption{Các tầng kiểm thử}\label{tab:tests}
\begin{tabularx}{\textwidth}{@{}L{3.3cm}L{1.6cm}Y@{}}
\toprule
\textbf{Tầng} & \textbf{Số lượng} & \textbf{Phạm vi} \\
\midrule
Python/PostgreSQL & 379 & Đợt hồi quy sau auth gồm đơn vị, tích hợp BR4/BR5 và hồi quy API BR6--BR8. \\
Spark/Delta trực tiếp & 2 & Snapshot khớp oracle với kết quả thiếu, sai, lỗi hạ tầng và nhiều phiên bản. Spark phục hồi checkpoint, khử trùng và rút lại hoạt động. \\
Đầu cuối CDC & 6 kiểm tra & Sửa điểm qua API, phát lại, UPDATE/DELETE hoạt động và phục hồi checkpoint. Phép thử riêng chèn lỗi publish xác nhận mart cũ vẫn phục vụ. \\
Trình duyệt & 5 & Sau cập nhật auth, Playwright đạt ở ba chế độ, gồm đăng nhập staff, Code Lab Classroom và trạng thái Replay dùng chung. Các lượt chấp nhận CDC và BR6--BR8 có nhật ký riêng. \\
\bottomrule
\end{tabularx}
\end{table}

Ngoài ra, mã Python qua \code{ruff} và \code{mypy}; kho dữ liệu chạy cả 10 cổng chất lượng sau mỗi lần build.

\subsection{Truy vết yêu cầu và kết quả sử dụng}

{\small
\begin{longtable}{@{}L{2.25cm}L{3.1cm}L{4.3cm}L{4.85cm}@{}}
\caption{Từ yêu cầu và dữ liệu đến kết quả, minh chứng và giới hạn}\label{tab:scenarios}\\
\toprule
\textbf{BR và người dùng} & \textbf{Dữ liệu} & \textbf{Kết quả sử dụng} & \textbf{Minh chứng và giới hạn} \\
\midrule
\endfirsthead
\toprule
\textbf{BR và người dùng} & \textbf{Dữ liệu} & \textbf{Kết quả sử dụng} & \textbf{Minh chứng và giới hạn} \\
\midrule
\endhead
BR1 Người học & Sự kiện bài tập, hạn nộp, học liệu & Xem thông báo phù hợp với khóa học & Giao diện demo và kiểm thử luật. Tích hợp từ dịch vụ Kafka đến hộp thông báo cần nghiệm thu. \\
BR2 Người học & Ghi danh, lịch sử học, đặc trưng dài hạn & Xem gợi ý kèm lý do & Luật gợi ý và bảng đặc trưng lô. Chưa có nhánh phiên qua feature store. \\
BR3 Người học & Mã nộp, stderr, diff, kết quả từng test & Chẩn đoán lỗi và nộp lại ở Code Lab & Classroom giữ cùng ID đến năm kết quả test và facts BR4. Luồng worker Kafka riêng chưa được nối đầu cuối với giao diện này. \\
BR4 Giảng viên & Bài nộp và kết quả chấm theo phiên bản & Khoanh vùng test bất thường và xem tác động sửa điểm & $2/24$ thành $3/24$ qua CDC. Oracle, kiểm tra bản công bố và lỗi publish. Bốn trường hợp gắn nhãn là fixture tổng hợp. \\
BR5 Cố vấn & Hoạt động, ghi danh, deadline và kết quả nộp & Ưu tiên can thiệp và theo dõi phục hồi & 100 xuống 50 điểm, một cảnh báo được giải quyết sau phát lại. OULAD đo khả năng dự báo hồi cứu. \\
BR6 Quản trị viên & Ghi danh và kết quả đã khử định danh & So sánh tiến độ tổ chức được phép & Acme và Globex minh họa ngưỡng $k=5$. Nhóm nhỏ ẩn tỷ lệ và số hoàn thành. Truy cập ngoài phạm vi trả 403. \\
BR7 Giảng viên & Tương tác theo giây, phiên bản video & Xác định đoạn cần xem lại & Hai phiên bản video có heatmap riêng. Chỉ số \en{seek} hiện gộp cả tua tới và tua lùi. \\
BR8 Giảng viên & Phiên bản transcript, CDC và sự kiện tương tác & Xuất bản, theo dõi tìm kiếm và xem xu hướng giảm & Các bài đo đồng bộ nội dung và lớp phủ ở Mục~\ref{chap:eval}. Embedding băm chủ yếu phản ánh từ vựng. \\
\bottomrule
\end{longtable}}

Trong bốn kịch bản BR4 có nhãn, oracle và phân loại khớp: test đáng ngờ, câu hỏi tỷ lệ thấp, đề bình thường và mẫu không đủ. Hai dương tính và hai âm tính đều đúng trên fixture này. Bài sửa điểm, bài đọc snapshot, bài restart cửa sổ và bài xác thực phạm vi là các minh họa kỹ thuật riêng. Phép thử Spark riêng kiểm tra bản ghi thiếu hoặc không hợp lệ, lỗi hạ tầng và phiên bản test. Khi chèn lỗi trong transaction publish, API giữ mart cũ. Diễn tập khôi phục so checksum các mart, gửi lại ID cũ, thêm hoạt động mới và xác nhận không trùng cảnh báo. Job kế thừa cần đợi lease hết hạn trước khi worker tiếp quản.

Trong diễn tập backfill, nhóm tách cơ sở dữ liệu, connector, topic, nhóm tiêu thụ và checkpoint khỏi đường cũ. Đối chiếu số dòng với nguồn và hai cửa sổ với oracle trước khi chuyển API cùng writer. Sau chuyển đổi, sự kiện kiểm tra chỉ hạ cánh ở đích mới. Giữ deployment cũ để quay lui. Lệnh retention mặc định trả bản xem trước và kiểm tra công việc hoặc dữ liệu chờ trước khi xóa.

\subsection{Môi trường và phương pháp đo}

Các nhật ký engine BR4/BR5 ngày 12/11/2026 dùng fixture tổng hợp và revision \code{3cfd70d}. Máy Linux x86-64 có 6 CPU logic và khoảng 18{,}4 GiB RAM. Cấu hình đối chiếu ghi 2 CPU và 2 GiB cho các engine. BR4 dùng PostgreSQL 16.15 và Spark 3.5.5. BR5 dùng Spark 3.5.1 và Flink 1.19.3, publish mỗi 5 giây và Flink checkpoint mỗi 10 giây. Mỗi cỡ có một lượt làm nóng và ba lượt đo. Mỗi tình huống đầu cuối có năm lượt đo. Kết quả chính đạt 36/36 lượt xử lý, 30/30 lượt đầu cuối và 2/2 kiểm tra restart. Bảng ghi trung vị và khoảng nhỏ nhất đến lớn nhất. Phép đo BR6--BR8 ngày 02/11 dùng VM Colima 6 CPU, 10 GB, Spark 3.5.5, Delta 3.3.2, Trino 478 và OpenSearch 2.17.1.

\subsection{Kết quả hiệu năng}

\subsubsection{BR4: PostgreSQL SQL và Spark SQL}

Đồng hồ chặng xử lý chạy từ lúc gọi tổng hợp đến khi ghi mart. Spark đọc cùng dữ liệu đã chuẩn hóa qua JDBC. Đường đo này gồm tổng hợp cửa sổ hiện tại và kỳ trước và ghi PostgreSQL. Đường triển khai Spark/Delta còn phải hạ cánh Delta, phân tích cohort và ghi facts, dimensions theo từng bản công bố. Kịch bản sửa điểm của harness đo từ sau commit nguồn đến khi API đọc được 3 lượt đạt, với harness điều phối lần tính.

\input{data/br45-engine-tables}
\input{data/br4-engine-chart}

Ở cả ba cỡ PostgreSQL có trung vị thấp hơn trong chặng đo BR4. Chi phí khởi động và điều phối Spark ảnh hưởng rõ khi cỡ nhỏ. Nhóm giữ PostgreSQL làm phương án phục vụ thay thế. Đường Spark/Delta hỗ trợ truy vết theo phiên bản và đọc lại snapshot cho điều tra lỗi chấm.

\subsubsection{BR4: sửa điểm qua worker triển khai}

Hai project Compose cô lập dùng cùng fixture 24 học viên, mốc 16:00 UTC ngày 12/11/2026 và cùng cấu hình CDC. Worker có giới hạn 2 CPU và 2.304 MiB. Hai đường dùng PostgreSQL 16.15. Đường Spark 3.5.5 và Delta 3.3.2 lưu snapshot và bản công bố trên MinIO. Nhật ký ngày 13/11 ghi bản mã API \code{eab34a8} kèm thay đổi auth, bản mã image Spark \code{1ac28e1}, image ID và giới hạn của từng dịch vụ.

Đồng hồ bắt đầu ngay sau commit sửa một lần chấm từ trượt sang đạt. Client thăm dò API mỗi 250 ms đến khi bản công bố thành công trả 3/24, đúng engine, có publication và source revision mới. Đường Spark phải có Delta version tăng. Mỗi engine có một lượt làm nóng và năm lượt đo. Sau mỗi lượt, client trả fixture về 2/24 và đợi công bố trước lần sửa tiếp theo. Bảng~\ref{tab:br4-deployed} tính độ trễ trên các lượt đạt, với giới hạn chờ 480 giây và nhật ký riêng cho từng lỗi.

\input{data/br45-deployed-table}

Phép đo này bao gồm PostgreSQL WAL, Debezium, Kafka, worker và lần đọc API. Đường Spark còn chờ hạ cánh snapshot Delta và ghi facts, dimensions trước khi thay mart. Trong các lượt đã quan sát, PostgreSQL công bố sớm hơn. Máy dùng chung có tải nền thay đổi, gồm thời gian khôi phục kho dữ liệu cho ảnh. Vì vậy độ chênh quan sát được gộp cả chi phí của đường xử lý và thời gian tranh chấp tài nguyên. Nhật ký giữ riêng lượt PostgreSQL bị chồng thao tác trong lúc phục hồi dịch vụ.

\subsubsection{BR5: Spark Structured Streaming và Flink}

Chặng xử lý bắt đầu khi gửi sự kiện vào Kafka và kết thúc khi sink PostgreSQL có đúng cửa sổ, đồng thời offsets xác nhận đã tiêu thụ. Thông lượng lấy số sự kiện duy nhất chia cho thời gian quan sát. Đầu vào còn có khoảng 1\% bản sao và 16 sự kiện đồng hồ ở mỗi lượt. Kiểm tra đúng đắn gồm trùng, sai thứ tự, biên nửa mở và rút đóng góp. Trong kịch bản cảnh báo và phục hồi, đồng hồ bắt đầu sau commit hoạt động nguồn và dừng khi API trả điểm và trạng thái cảnh báo đúng. Hai engine dùng chung bộ tính và lưu điểm sau sink benchmark. Cả hai khôi phục đúng tổng cửa sổ và định danh sau restart.

\input{data/br5-engine-tables}
\input{data/br5-engine-chart}

Ở ba cỡ đã đo, trung vị chặng Kafka đến cửa sổ của Spark thấp hơn Flink trong cấu hình này. Một lượt nghiệm thu riêng trên đường Spark triển khai ghi nhận 20{,}97 giây từ CDC đến cảnh báo và 35{,}46 giây từ restart đến phục hồi. Hai thời gian nghiệm thu mô tả đường chạy thực tế và không được gộp vào phân phối benchmark.

\paragraph{Số liệu lịch sử tháng 10.} Ở 100.000 lượt chấm, trung vị BR4 là 3{,}90 giây cho PostgreSQL SQL và 1{,}95 giây cho tính toán Python trên dữ liệu đã chuẩn hóa. Ở 100.000 sự kiện, trung vị BR5 là 15{,}17 giây từ Kafka đã chuẩn hóa đến khi lưu cửa sổ Spark. Các số này dùng ranh giới của phép đo cũ. CSV gốc nằm trong \code{data/}.

\subsubsection{BR6--BR8: phục vụ mart}

\begin{figure}[H]
\centering
\begin{tikzpicture}
\begin{axis}[
  ybar, bar width=16pt, width=0.82\textwidth, height=6.2cm,
  ymode=log, log basis y=10, ymin=0.8, ymax=400,
  symbolic x coords={BR6,BR7,BR8}, xtick=data,
  xticklabels={{BR6: tiến độ tổ chức}, {BR7: heatmap 390 s}, {BR8: tương tác 38 học viên}},
  ylabel={Trung vị độ trễ (ms, trục log)},
  ymajorgrids, enlarge x limits=0.25,
  legend style={at={(0.5,1.02)}, anchor=south, legend columns=2},
  nodes near coords, nodes near coords style={font=\scriptsize, color=ink},
  point meta=rawy,
]
\addplot[fill=sone, draw=none] coordinates {(BR6,63.0) (BR7,92.1) (BR8,41.6)};
\addlegendentry{Đọc kho qua Trino, chưa cache}
\addplot[fill=stwo, draw=none] coordinates {(BR6,1.4) (BR7,2.5) (BR8,1.6)};
\addlegendentry{API, cache mart đã ấm}
\end{axis}
\end{tikzpicture}
\caption{Độ trễ phục vụ các mart BR6--BR8 (30 request mỗi trường hợp)}
\label{fig:serving}
\end{figure}

Bách phân vị 90 khi đọc thẳng qua Trino là 92, 136 và 59 ms cho BR6, BR7, BR8; khi có cache là 1,8, 2,8 và 2,0 ms. Vì mart chỉ thay đổi mỗi ngày còn dashboard thăm dò mỗi 5 giây, bộ đọc có cache TTL 60 giây (\code{CachedMartReader}). Cache này còn sửa được một lỗi thật: Trino bị hết bộ nhớ khi dashboard thăm dò liên tục mà không có cache.

\subsubsection{BR8: từ xuất bản đến khi tìm kiếm được}

\begin{figure}[H]
\centering
\begin{tikzpicture}
\begin{axis}[
  width=0.82\textwidth, height=5.6cm,
  xlabel={Lần xuất bản}, ylabel={Độ trễ (ms)},
  xmin=0.5, xmax=20.5, ymin=0, ymax=1200,
  xtick={1,5,10,15,20}, ymajorgrids,
]
\addplot[color=sone, mark=*, mark size=2pt, line width=0.9pt]
  table[x=publish, y=publish_to_searchable_ms, col sep=comma] {data/content_sync_latency.csv};
\draw[stgood!70!black, dashed] (axis cs:0.5,535) -- (axis cs:20.5,535)
  node[pos=0.98, above, anchor=south east, font=\scriptsize, color=ink]{trung vị 535 ms};
\end{axis}
\end{tikzpicture}
\caption{BR8: độ trễ từ xuất bản đến khi tìm kiếm được, 20 lần xuất bản liên tiếp}
\label{fig:sync-latency}
\end{figure}

API nhận yêu cầu xuất bản trong 16,9 ms (trung vị). Toàn bộ đường đi PostgreSQL WAL $\rightarrow$ Debezium $\rightarrow$ Kafka $\rightarrow$ content-sync $\rightarrow$ tìm kiếm được có trung vị \textbf{535 ms}, p90 992 ms, lớn nhất 1.024 ms. Bản thân bước đồng bộ chỉ mất 25--70 ms; phần còn lại là thời gian bắt thay đổi và chu kỳ thăm dò của consumer.

\subsubsection{BR7/BR8: lớp phủ thời gian thực}

\begin{table}[H]
\centering\small
\caption{Độ trễ từ sự kiện đến khi hiển thị trong API (3 đợt mỗi kích thước)}\label{tab:live}
\begin{tabular}{@{}rrrc@{}}
\toprule
\textbf{Phiên mỗi đợt} & \textbf{Sự kiện thô (trung bình)} & \textbf{Trung bình (nhỏ nhất -- lớn nhất)} & \textbf{Chính xác} \\
\midrule
25 & 234 & 7,27 s (7,19 -- 7,35) & có \\
100 & 888 & 7,23 s (6,95 -- 7,38) & có \\
400 & 3.622 & 7,12 s (5,69 -- 8,03) & có \\
\bottomrule
\end{tabular}
\end{table}

Độ trễ gần như không đổi khi tải tăng 15 lần, vì nó bị chi phối bởi chi phí cố định của mỗi micro-batch (bước chuẩn hoá khoảng 2,7 s, trigger 1 s) chứ không phải khối lượng. ``Chính xác'' nghĩa là lớp phủ chứa đúng số tương tác đúng hạn kỳ vọng của đợt cho video nhiều tương tác nhất.

\subsubsection{Thời gian build kho dữ liệu}

\begin{figure}[H]
\centering
\begin{tikzpicture}
\begin{axis}[
  ybar stacked, bar width=26pt, width=0.8\textwidth, height=6.4cm,
  symbolic x coords={35,70,140}, xtick=data,
  xticklabels={{35 ngày\\10.761 sự kiện}, {70 ngày\\20.486 sự kiện}, {140 ngày\\40.221 sự kiện}},
  xticklabel style={align=center},
  ylabel={Thời gian (s)}, ymin=0, ymax=160, ymajorgrids, enlarge x limits=0.25,
  legend style={at={(0.5,1.03)}, anchor=south, legend columns=2, legend cell align=left},
]
\addplot[fill=sone, draw=white] coordinates {(35,32.6) (70,36.1) (140,37.3)};
\addlegendentry{silver}
\addplot[fill=stwo, draw=white] coordinates {(35,54.9) (70,48.8) (140,51.1)};
\addlegendentry{gold}
\addplot[fill=sthree, draw=white] coordinates {(35,11.6) (70,6.3) (140,7.0)};
\addlegendentry{mart}
\addplot[fill=sfour, draw=white] coordinates {(35,44.6) (70,37.4) (140,40.4)};
\addlegendentry{khởi tạo, kiểm tra chất lượng, khác}
\end{axis}
\end{tikzpicture}
\caption{Thời gian build toàn bộ 43 model theo khối lượng dữ liệu (trung bình 2 lần)}
\label{fig:build}
\end{figure}

Tổng thời gian build là 143,7, 128,6 và 135,8 giây cho 35, 70 và 140 ngày dữ liệu; mọi model đều thành công. Thời gian gần như phẳng khi khối lượng tăng 4 lần: ở quy mô này, chi phí cố định của mỗi model (khởi tạo job Spark, commit Delta, kiểm tra chất lượng) chiếm phần lớn. Các mart BR6--BR8 chỉ mất 1--6 giây. Chênh lệch giữa hai lần chạy (123 và 164 giây ở 35 ngày) là nhiễu do máy ảo dùng chung với phần còn lại của stack.

\subsection{Hiệu quả cảnh báo BR5 trên OULAD}

Dữ liệu công khai OULAD~\cite{kuzilek2017oulad,ucioulad} có 10.655.280 dòng tương tác tổng hợp theo ngày trong lần chạy này. Phép đánh giá tạo 77.106 quan sát tại các ngày 14, 28 và 56. Đặc trưng có trước mốc dự báo, nhãn là rút học trong 30 ngày kế tiếp. Nhóm tách học viên 2013 bằng băm để hiệu chỉnh và xác thực. Tập 2014 loại học viên đã xuất hiện năm 2013 và được trình bày như đánh giá hồi cứu vì đã được xem trong các lượt trước.

\input{data/br5-oulad-table}

Ngưỡng 65 hòa ngưỡng 60 trên tập xác thực 2013 với F1 khoảng 0{,}128. Nhóm giữ policy v1 ở ngưỡng 60. PR-AUC trong bảng là average precision, còn thời gian báo trước là trung vị của dương tính đúng. Baseline không hoạt động bảy ngày báo sớm hơn ở tập xác thực, với F1 thấp hơn. Cố vấn cần dùng lời giải thích và lịch sử để kiểm tra các cảnh báo vì precision còn thấp. Các mốc lặp của một học viên có tương quan. Tổng hợp theo ngày cũng giới hạn khả năng suy luận về hoạt động trong ngày. Kết quả chưa đo hiệu quả can thiệp.

\subsection{Thảo luận}

Kết quả BR4 cho phép giảng viên ưu tiên kiểm tra test lỗi và xác minh việc sửa điểm qua publication mới. Với BR5, cố vấn thấy các tín hiệu đóng góp, trạng thái xác nhận, ghi chú và phục hồi trên cùng cảnh báo. Độ trễ ở đây phụ thuộc chu kỳ publish, thời gian truy vấn và tài nguyên của máy đo. Lựa chọn engine cần được đánh giá lại khi thay đổi tải hoặc cấu hình. Các minh chứng giao diện ở Mục~\ref{chap:ui} nối kết quả kỹ thuật với thao tác nghiệp vụ.
